% =====================================================================
%  LIBRO ACADÉMICO: APLICACIONES WEB - EDICIÓN AMPLIADA
%  CAPÍTULO 4 (SEMANA 4): JavaScript moderno y programación asíncrona
%  Dr. Gleiston Cicerón Guerrero Ulloa, Ph.D. — UTEQ — 2026
% =====================================================================
\documentclass[11pt,a4paper,openany]{book}

\usepackage[utf8]{inputenc}
\usepackage[T1]{fontenc}
\usepackage[spanish,es-tabla,es-noshorthands]{babel}
\usepackage{lmodern}
\usepackage{microtype}
\usepackage[a4paper,margin=2.5cm,headheight=14.5pt]{geometry}
\usepackage{setspace}\onehalfspacing
\usepackage{parskip}

\usepackage{xcolor}
\definecolor{uteqverde}{RGB}{0,102,51}
\definecolor{uteqdorado}{RGB}{204,153,0}
\definecolor{codebg}{RGB}{248,248,242}
\definecolor{codeborder}{RGB}{180,180,180}
\definecolor{codekey}{RGB}{0,87,174}
\definecolor{codestr}{RGB}{196,26,22}
\definecolor{codecom}{RGB}{113,113,113}
\definecolor{codefn}{RGB}{123,47,190}
\definecolor{codeobj}{RGB}{11,107,138}
\definecolor{codeop}{RGB}{200,123,0}
\definecolor{codeident}{RGB}{45,128,0}
\definecolor{codenum}{RGB}{56,58,66}

\usepackage{amsmath,amssymb}
\usepackage{graphicx}
\usepackage{float}
\usepackage{tikz}
\usetikzlibrary{shapes,arrows,arrows.meta,positioning,fit,backgrounds,calc}
\usetikzlibrary{shapes.geometric,shadows,matrix,decorations.markings}
\usetikzlibrary{babel}

\usepackage{fancyhdr}
\pagestyle{fancy}\fancyhf{}
\fancyhead[LE,RO]{\thepage}
\fancyhead[RE]{\nouppercase{\leftmark}}
\fancyhead[LO]{\nouppercase{\rightmark}}
\renewcommand{\headrulewidth}{0.4pt}

\usepackage[most]{tcolorbox}
\tcbuselibrary{breakable,skins,listings}

\newtcolorbox{cajadefinicion}[1]{colback=uteqverde!5,colframe=uteqverde,
	fonttitle=\bfseries\color{white},title=#1,breakable,enhanced}
\newtcolorbox{cajaejemplo}[1]{colback=uteqdorado!8,colframe=uteqdorado!80!black,
	fonttitle=\bfseries\color{white},title=#1,breakable,enhanced}
\newtcolorbox{cajaejercicio}[1]{colback=blue!4,colframe=blue!60!black,
	fonttitle=\bfseries\color{white},title=#1,breakable,enhanced}
\newtcolorbox{cajaadvertencia}[1]{colback=red!5,colframe=red!70!black,
	fonttitle=\bfseries\color{white},title=#1,breakable,enhanced}
\newtcolorbox{cajahistoria}[1]{colback=blue!3,colframe=blue!50!black,
	fonttitle=\bfseries\color{white},title=#1,breakable,enhanced}
\newtcolorbox{cajabuenas}[1]{colback=green!4,colframe=uteqverde!90!black,
	fonttitle=\bfseries\color{white},title=#1,breakable,enhanced}
\newtcolorbox{cajacompara}[1]{colback=cyan!5,colframe=cyan!60!black,
	fonttitle=\bfseries\color{white},title=#1,breakable,enhanced}

% ---------- Listings ---------------------------------------------------
\usepackage{listings}

\lstdefinelanguage{JavaScript}{
	keywords={break,case,catch,continue,debugger,default,delete,do,else,
		finally,for,function,if,in,instanceof,new,return,switch,this,throw,
		try,typeof,var,void,while,with,let,const,of,yield,async,await,
		class,extends,super,static,import,export,from,as,true,false,null,
		undefined,get,set},
	ndkeywords={Array,Boolean,Date,Error,Function,JSON,Math,Number,Object,
		Promise,RegExp,String,Symbol,console,document,window,fetch,Map,Set,
		WeakMap,WeakSet,Proxy,Reflect,setTimeout,setInterval,clearTimeout,
		clearInterval,parseInt,parseFloat,isNaN,isFinite,encodeURIComponent,
		decodeURIComponent,localStorage,sessionStorage,navigator,history,
		location,XMLHttpRequest,AbortController,URL,URLSearchParams,
		MutationObserver,IntersectionObserver,EventSource,Worker},
	sensitive=true,
	comment=[l]{//},
	morecomment=[s]{/*}{*/},
	morestring=[b]',
	morestring=[b]",
	morestring=[b]`
}

\lstdefinestyle{codigobase}{
	basicstyle=\ttfamily\footnotesize\color{codenum},
	backgroundcolor=\color{codebg},
	frame=single,
	rulecolor=\color{codeborder},
	numbers=left,
	numberstyle=\tiny\color{gray!70},
	numbersep=8pt,
	stepnumber=1,
	breaklines=true,
	breakatwhitespace=true,
	showstringspaces=false,
	tabsize=2,
	keywordstyle=\color{codekey}\bfseries,
	ndkeywordstyle=\color{codeobj},
	stringstyle=\color{codestr},
	commentstyle=\color{codecom}\itshape,
	captionpos=b,
	xleftmargin=12pt,
	xrightmargin=4pt,
	inputencoding=utf8,
	extendedchars=true,
	literate=
	{á}{{\'a}}1 {é}{{\'e}}1 {í}{{\'i}}1 {ó}{{\'o}}1 {ú}{{\'u}}1
	{Á}{{\'A}}1 {É}{{\'E}}1 {Í}{{\'I}}1 {Ó}{{\'O}}1 {Ú}{{\'U}}1
	{ñ}{{\~n}}1 {Ñ}{{\~N}}1 {¿}{{?`}}1 {¡}{{!`}}1
	{—}{{---}}1 {…}{{...}}1
}
\lstdefinestyle{js}{style=codigobase,language=JavaScript}
\lstdefinestyle{html}{style=codigobase,language=HTML}
\lstdefinestyle{shell}{style=codigobase,language=bash}

% ---------- Tablas -----------------------------------------------------
\usepackage{array}\usepackage{booktabs}\usepackage{longtable}
\usepackage{tabularx}\usepackage{multirow}

% ---------- Hiperenlaces -----------------------------------------------
\usepackage[hidelinks,bookmarks=true,bookmarksnumbered=true]{hyperref}
\hypersetup{colorlinks=true,linkcolor=uteqverde,citecolor=uteqverde,
	urlcolor=uteqdorado!70!black,
	pdftitle={Cap. 4 JavaScript Moderno - Aplicaciones Web UTEQ},
	pdfauthor={Dr. Gleiston Ciceron Guerrero Ulloa}}

% ---------- Listas -----------------------------------------------------
\usepackage{enumitem}\setlist{itemsep=2pt,topsep=2pt}

% ---------- Epígrafe ---------------------------------------------------
\usepackage{epigraph}\setlength{\epigraphwidth}{0.72\textwidth}

% ---------- Formato de títulos -----------------------------------------
\usepackage{titlesec}
\titleformat{\chapter}[display]
{\normalfont\Huge\bfseries\color{uteqverde}}
{\filleft\Large Semana \thechapter}
{1ex}{\titlerule\vspace{1ex}\filright}[\vspace{0.5ex}\titlerule]
\titleformat{\section}
{\normalfont\Large\bfseries\color{uteqverde}}{\thesection}{1em}{}
\titleformat{\subsection}
{\normalfont\large\bfseries\color{uteqverde!85!black}}{\thesubsection}{1em}{}
\titleformat{\subsubsection}
{\normalfont\normalsize\bfseries\color{uteqdorado!90!black}}{\thesubsubsection}{1em}{}

% ---------- Contador de capítulo ---------------------------------------
\setcounter{chapter}{3}   % el siguiente \chapter será el 4

% =====================================================================
\begin{document}
	
	\setlength{\hfuzz}{3pt}
	
	% --------------------- Portada del capítulo --------------------------
	\begin{titlepage}
		\centering\vspace*{1.5cm}
		{\color{uteqverde}\rule{\linewidth}{2pt}}\\[0.8cm]
		{\Huge\bfseries\color{uteqverde} APLICACIONES WEB}\\[0.5cm]
		{\Large\itshape Edición ampliada — Historia, fundamentos, prácticas y
			buenas prácticas profesionales}\\[0.3cm]
		{\color{uteqverde}\rule{\linewidth}{1pt}}\\[1.5cm]
		{\LARGE\bfseries Semana 4}\\[0.5cm]
		{\Huge\bfseries\color{uteqdorado!80!black}
			JavaScript moderno\\[0.3cm]y programación asíncrona}\\[2cm]
		{\large\bfseries Universidad Técnica Estatal de Quevedo}\\
		{\large Facultad de Ciencias de la Computación}\\
		{\large Carrera de Ingeniería de Software (Rediseño)}\\[0.7cm]
		{\large\itshape Quinto nivel — Período Académico 2026--2027 PPA}\\[1.5cm]
		{\large \textbf{Autor:} Dr.\ Gleiston Cicerón Guerrero Ulloa, Ph.D.}\\[2cm]
		{\color{uteqdorado}\rule{\linewidth}{2pt}}\\[0.5cm]
		{\large Quevedo, Los Ríos — Ecuador \hfill Mayo de 2026}
	\end{titlepage}
	
	\tableofcontents
	\newpage
	
	% =====================================================================
	%   CAPÍTULO 4 — JavaScript moderno y programación asíncrona
	% =====================================================================
	\chapter{JavaScript moderno y programación asíncrona}
	\label{cap:semana4}
	
	\epigraph{«Cualquier aplicación que pueda escribirse en JavaScript,
		eventualmente se escribirá en JavaScript.»}
	{--- \textit{Jeff Atwood, 2007 (Ley de Atwood)}}
	
	% ---------- Resultado de aprendizaje ----------------------------------
	\section*{Resultado de aprendizaje de la semana}
	\addcontentsline{toc}{section}{Resultado de aprendizaje de la semana}
	
	Al concluir la semana~4, el estudiante:
	\begin{itemize}
		\item Comprende la evolución histórica y el modelo de ejecución de
		JavaScript en el navegador.
		\item Declara variables con \texttt{let}/\texttt{const}, escribe
		funciones flecha, aplica \emph{destructuring} y \emph{spread}.
		\item Lee y modifica el árbol DOM, escucha eventos con el patrón
		\texttt{addEventListener} y aplica delegación de eventos.
		\item Comprende el \emph{event loop}, la pila de llamadas y la cola
		de tareas.
		\item Escribe funciones de orden superior con \texttt{.map()},
		\texttt{.filter()} y \texttt{.reduce()}.
		\item Modela y reutiliza comportamiento con cierres (\emph{closures})
		y con clases ES6.
		\item Consume APIs REST con \texttt{fetch} y \texttt{async/await},
		encadena Promises y gestiona errores con \texttt{try/catch}.
		\item Organiza el código en módulos ES con \texttt{import}/\texttt{export}.
		\item Aplica las doce buenas prácticas profesionales de JavaScript
		2026 \cite{flanagan2020javascript,crockford2008javascript}.
	\end{itemize}
	
	% =====================================================================
	\section{Historia de JavaScript}
	\label{sec:4.historia}
	
	JavaScript nació en 1995 como un lenguaje «de juguete» creado en
	diez días. Hoy es la \emph{lingua franca} del navegador, ejecuta
	servidores enteros (Node.js), corre en satélites y dispositivos IoT,
	y es el lenguaje más usado del mundo según el índice TIOBE 2025.
	Conocer su historia ayuda a entender por qué tiene particularidades
	sintácticas y por qué evoluciona tan rápido.
	
	\begin{cajahistoria}{De diez días de diseño a la \emph{lingua franca} de la web}
		En mayo de 1995, Netscape contrató a \textbf{Brendan Eich} para
		crear un lenguaje de scripting que se ejecutara en el navegador.
		Le dieron \emph{diez días}. El resultado, llamado originalmente
		\textbf{Mocha}, luego \textbf{LiveScript} y finalmente
		\textbf{JavaScript} (1996, por razones de marketing vinculadas a la
		popularidad de Java), reunía rasgos memorables: tipado dinámico,
		\emph{first-class functions}, prototipos en lugar de clases,
		\texttt{undefined} y \texttt{null} como dos formas distintas de
		«nada» \cite{crockford2008javascript}.
		
		Microsoft respondió con JScript en Internet Explorer 3 (1996),
		provocando una guerra de incompatibilidades que se cerró con la
		estandarización ECMA en 1997 (\textbf{ECMAScript~1}). Durante una
		década, JavaScript fue considerado un lenguaje de segunda clase.
		El cambio llegó con el artículo de Jesse James Garrett que acuñó
		\textbf{Ajax} (2005) \cite{garrett2005ajax} y con el motor V8 de
		\textbf{Google Chrome} (2008), que compilaba JS a código nativo y
		era hasta 100× más rápido que sus predecesores. Ryan Dahl aprovechó
		V8 para crear \textbf{Node.js} (2009), llevando JS al servidor.
		
		\textbf{ECMAScript 2015 (ES6)} fue el mayor salto del lenguaje:
		clases, módulos, \texttt{let}/\texttt{const}, funciones flecha,
		Promises y \emph{destructuring}. Desde entonces el estándar se
		actualiza cada año. ECMAScript 2024 (ES15) es hoy el estándar
		vigente \cite{ecma2024}; ES2026 incorpora
		\texttt{Array.fromAsync}, \texttt{Promise.try} y
		\emph{records/tuples} inmutables.
	\end{cajahistoria}
	
	\subsection{Tabla evolutiva de ECMAScript}
	\label{sub:4.ecma}
	
	\begin{table}[htbp]
		\centering
		\caption{Versiones más relevantes de ECMAScript.}
		\label{tab:4.ecma}
		\footnotesize
		\begin{tabularx}{\textwidth}{llX}
			\toprule
			\textbf{Versión} & \textbf{Año} & \textbf{Aportes principales}\\
			\midrule
			ES3        & 1999 & Expresiones regulares, \texttt{try/catch}, manejo de excepciones.\\
			ES5        & 2009 & \texttt{strict mode}, JSON nativo, \texttt{Array.forEach}, getters/setters.\\
			ES6 (ES2015) & 2015 & \texttt{let}/\texttt{const}, clases, módulos, arrow functions, Promises,
			template literals, \emph{destructuring}, \emph{spread}, \texttt{Map}/\texttt{Set}.\\
			ES2016     & 2016 & \texttt{Array.includes()}, operador exponente \texttt{**}.\\
			ES2017     & 2017 & \texttt{async/await}, \texttt{Object.entries()}, \texttt{String.padStart()}.\\
			ES2018     & 2018 & \emph{Rest/spread} en objetos, \texttt{Promise.finally()}, iteradores asíncronos.\\
			ES2019     & 2019 & \texttt{Array.flat()}, \texttt{Object.fromEntries()}, \texttt{String.trimStart()}.\\
			ES2020     & 2020 & Optional chaining (\texttt{?.}), nullish coalescing (\texttt{??}), BigInt.\\
			ES2021     & 2021 & \texttt{String.replaceAll()}, \texttt{Promise.any()}, \texttt{WeakRef}.\\
			ES2022     & 2022 & Top-level \texttt{await}, campos privados (\texttt{\#}), \texttt{Array.at()}.\\
			ES2023     & 2023 & \texttt{Array.findLast()}, \texttt{Array.toSorted()}, \texttt{Array.toReversed()}.\\
			ES2024     & 2024 & \texttt{Object.groupBy()}, \texttt{ArrayBuffer} transferible,
			\texttt{Promise.withResolvers()}.\\
			ES2026 \cite{ecma2024} & 2026 & \texttt{Array.fromAsync()}, \texttt{Promise.try()},
			\emph{records} y \emph{tuples} inmutables.\\
			\bottomrule
		\end{tabularx}
	\end{table}
	
	% =====================================================================
	\section{El modelo de ejecución: event loop y call stack}
	\label{sec:4.eventloop}
	
	Antes de escribir código asíncrono es indispensable comprender
	\emph{por qué} JavaScript necesita asincronía. El lenguaje corre
	en un único hilo (\emph{single-threaded}): solo ejecuta una
	instrucción a la vez. Si una operación lenta —una petición HTTP,
	una lectura de disco— bloqueara el hilo, la página quedaría
	congelada. El \textbf{event loop} resuelve este problema delegando
	las operaciones lentas al entorno (navegador o Node.js) y
	procesando sus resultados cuando el hilo principal está libre.
	
	\begin{cajadefinicion}{Glosario del modelo de ejecución}
		\begin{description}[leftmargin=4em,style=nextline]
			\item[\textbf{Call Stack}] Pila de llamadas donde se apilan los
			marcos de ejecución de las funciones activas. LIFO (último en
			entrar, primero en salir).
			\item[\textbf{Web APIs}] Servicios del navegador que gestionan
			operaciones asíncronas: \texttt{setTimeout}, \texttt{fetch},
			\texttt{XMLHttpRequest}, eventos del DOM, etc.
			\item[\textbf{Callback / Task Queue}] Cola donde las Web APIs
			depositan los callbacks listos para ejecutarse.
			\item[\textbf{Microtask Queue}] Cola de mayor prioridad para
			callbacks de Promises (\texttt{.then()}, \texttt{await}) y
			\texttt{MutationObserver}.
			\item[\textbf{Event Loop}] Bucle que comprueba continuamente si
			la Call Stack está vacía; si lo está, mueve primero todas las
			microtareas y luego una tarea de la Task Queue.
		\end{description}
	\end{cajadefinicion}
	
	La Figura~\ref{fig:4.eventloop} muestra el modelo completo.
	
	\begin{figure}[htbp]
		\centering
		\begin{tikzpicture}[
			box/.style={rectangle,draw=uteqverde,fill=uteqverde!8,thick,
				rounded corners=4pt,minimum width=3.2cm,minimum height=1cm,
				font=\small\bfseries,align=center},
			qbox/.style={rectangle,draw=uteqdorado!70!black,fill=uteqdorado!10,thick,
				rounded corners=4pt,minimum width=3.6cm,minimum height=0.85cm,
				font=\small,align=center},
			arr/.style={-{Latex[length=2.5mm]},thick,uteqverde},
			arrd/.style={-{Latex[length=2.5mm]},thick,uteqdorado!80!black},
			lbl/.style={font=\scriptsize,text=gray!70!black}
			]
			% Call Stack
			\node[box,minimum height=3.8cm,minimum width=3.2cm,align=center]
			(stack) at (0,0) {};
			\node[above=0.05cm of stack,font=\small\bfseries\color{uteqverde}]
			{Call Stack};
			\node[rectangle,draw=uteqverde!60,fill=white,minimum width=2.8cm,
			minimum height=0.6cm,font=\scriptsize] at (0,0.9) {frame 2: \texttt{render()}};
			\node[rectangle,draw=uteqverde!60,fill=white,minimum width=2.8cm,
			minimum height=0.6cm,font=\scriptsize] at (0,0.2) {frame 1: \texttt{main()}};
			\node[rectangle,draw=uteqverde!60,fill=uteqverde!5,minimum width=2.8cm,
			minimum height=0.6cm,font=\scriptsize\itshape] at (0,-0.5) {(vacío)};
			
			% Web APIs
			\node[box,minimum height=3.2cm,minimum width=3.4cm]
			(webapis) at (5.5,0) {};
			\node[above=0.05cm of webapis,font=\small\bfseries\color{uteqverde}]
			{Web APIs};
			\node[rectangle,draw=uteqverde!40,fill=white,minimum width=3cm,
			minimum height=0.52cm,font=\scriptsize] at (5.5,0.8) {\texttt{setTimeout(fn,500)}};
			\node[rectangle,draw=uteqverde!40,fill=white,minimum width=3cm,
			minimum height=0.52cm,font=\scriptsize] at (5.5,0.2) {\texttt{fetch(url)}};
			\node[rectangle,draw=uteqverde!40,fill=white,minimum width=3cm,
			minimum height=0.52cm,font=\scriptsize] at (5.5,-0.4) {\texttt{addEventListener}};
			
			% Task Queue
			\node[qbox,minimum width=3.6cm,minimum height=0.9cm]
			(taskq) at (5.5,-2.8) {\textbf{Task Queue} (Macrotareas)\\
				\scriptsize click, timeout, XHR…};
			
			% Microtask Queue
			\node[qbox,minimum width=3.6cm,minimum height=0.9cm,
			draw=cyan!50!black,fill=cyan!6]
			(microtaskq) at (0,-2.8) {\textbf{Microtask Queue}\\
				\scriptsize Promise.then / await};
			
			% Event Loop
			\node[circle,draw=uteqdorado!80!black,fill=uteqdorado!15,thick,
			minimum size=1.5cm,font=\scriptsize\bfseries,align=center]
			(loop) at (2.75,-2.8) {Event\\Loop};
			
			% Arrows
			\draw[arr] (stack.east) -- node[above,lbl]{delega}(webapis.west);
			\draw[arrd] (webapis.south) -- ++(0,-0.5) -| (taskq.north);
			\draw[arrd] (taskq.west) -- (loop.east);
			\draw[arrd,cyan!70!black] (microtaskq.east) -- (loop.west);
			\draw[arr] (loop.north) -- ++(0,0.5) -| node[right,lbl,pos=0.3]
			{push a stack} (stack.south);
			\node[lbl,align=center] at (2.75,-4.0)
			{\textit{1º vacía microtask queue — 2º una tarea de task queue}};
		\end{tikzpicture}
		\caption{Modelo de ejecución de JavaScript: call stack, Web APIs,
			microtask queue, task queue y event loop.}
		\label{fig:4.eventloop}
	\end{figure}
	
	\begin{cajaejemplo}{Ejemplo 4.0 — Orden de ejecución del event loop}
		El siguiente código ilustra la prioridad de las colas.
		
		\begin{lstlisting}[style=js,numbers=none,caption={Orden observable en la consola},
			label=lst:4.0]
			console.log("1 - sincrono");                        // stack inmediato
			
			setTimeout(() => console.log("4 - macrotarea"), 0); // task queue
			
			Promise.resolve()
			.then(() => console.log("3 - microtarea"));       // microtask queue
			
			console.log("2 - sincrono");                        // stack inmediato
		\end{lstlisting}
		
		\textbf{Salida en consola:}
		\begin{lstlisting}[style=shell,numbers=none]
			1 - sincrono
			2 - sincrono
			3 - microtarea
			4 - macrotarea
		\end{lstlisting}
		
		\textbf{Explicación:} Las dos líneas síncronas se ejecutan primero
		(stack). Luego se vacía la microtask queue (la Promise). Solo
		entonces el event loop procesa la macrotarea del \texttt{setTimeout},
		aunque su retardo sea cero.
	\end{cajaejemplo}
	
	% =====================================================================
	\section{Fundamentos modernos del lenguaje}
	\label{sec:4.fundamentos}
	
	Esta sección cubre las características del lenguaje que todo
	desarrollador JavaScript 2026 debe dominar antes de abordar el DOM
	o el código asíncrono.
	
	% ----- 4.3.1 Variables -----------------------------------------------
	\subsection{\texttt{let}, \texttt{const} y el fin de \texttt{var}}
	\label{sub:4.letconst}
	
	La palabra clave \texttt{var} tenía dos defectos graves:
	\emph{hoisting} (las variables se elevaban silenciosamente al
	inicio del \emph{scope}) y \emph{function scope} en lugar de
	\emph{block scope}. ES6 los corrigió con \texttt{let}
	(mutable, block-scoped) y \texttt{const} (inmutable en
	referencia, block-scoped). Véase el Listado~\ref{lst:4.1}.
	
	\begin{lstlisting}[style=js,numbers=none,caption={Variables en JavaScript moderno},
		label=lst:4.1]
		// REGLA PRACTICA: const por defecto, let si muta, var NUNCA
		const PI = 3.14159;          // inmutable; error si se reasigna
		let   contador = 0;          // mutable, scope de bloque
		var   legacy   = "evitar";   // hoisted, scope de funcion - NO usar
		
		// Hoisting de var (trampa clasica):
		console.log(legacy2);        // undefined, NO un ReferenceError
		var legacy2 = "peligroso";
		
		// Block scope de let/const:
		{
			let local = 10;
			const MAX = 100;
		}
		// console.log(local); // ReferenceError: local is not defined
	\end{lstlisting}
	
	% ----- 4.3.2 Funciones flecha -----------------------------------------
	\subsection{Funciones regulares y funciones flecha (\texttt{=>})}
	\label{sub:4.arrow}
	
	Las funciones flecha son más concisas y, crucialmente, \textbf{no
		enlazan su propio \texttt{this}}: heredan el \texttt{this} del
	contexto léxico donde fueron definidas. Esto las hace ideales como
	callbacks pero inadecuadas como métodos de objetos cuando se
	necesita acceso a la instancia. Véase el Listado~\ref{lst:4.2}.
	
	\begin{lstlisting}[style=js,caption={Funciones regulares vs.\ funciones flecha},
		label=lst:4.2]
		// Funcion regular con declaracion
		function sumar(a, b) {
			return a + b;
		}
		
		// Funcion flecha equivalente (cuerpo conciso)
		const sumar = (a, b) => a + b;
		
		// Arrow sin parametros
		const ahora = () => new Date().toISOString();
		
		// Arrow con cuerpo completo
		const dividir = (a, b) => {
			if (b === 0) throw new Error("Division por cero");
			return a / b;
		};
		
		// Diferencia critica: this lexico
		function Temporizador() {
			this.segundos = 0;
			// INCORRECTO con funcion regular: this sera undefined en strict
			// setInterval(function() { this.segundos++; }, 1000);
			
			// CORRECTO con arrow: hereda this del constructor
			setInterval(() => {
				this.segundos++;
				console.log(this.segundos);
			}, 1000);
		}
		const t = new Temporizador();
	\end{lstlisting}
	
	% ----- 4.3.3 Template literals ----------------------------------------
	\subsection{Template literals y etiquetas}
	\label{sub:4.template}
	
	Los \emph{template literals} (backticks) permiten interpolación
	de expresiones, cadenas multilínea y \emph{tagged templates} para
	sanitización o internacionalización. Véase el Listado~\ref{lst:4.3}.
	
	\begin{lstlisting}[style=js,numbers=none,caption={Template literals},label=lst:4.3]
		const nombre = "Maria";
		const edad   = 22;
		
		// Concatenacion clasica (evitar)
		const msg1 = "Hola " + nombre + ", tienes " + edad + " anios.";
		
		// Template literal (preferido)
		const msg2 = `Hola ${nombre}, tienes ${edad} anios.`;
		
		// Expresion compleja dentro de ${}
		const precio = 9.5;
		const iva = 0.15;
		console.log(`Total con IVA: $${(precio * (1 + iva)).toFixed(2)}`);
		
		// Cadena multilinea
		const html = `
		<div class="tarjeta">
		<h2>${nombre}</h2>
		<p>Edad: ${edad}</p>
		</div>
		`;
	\end{lstlisting}
	
	% ----- 4.3.4 Destructuring y spread -----------------------------------
	\subsection{Destructuring, spread y rest}
	\label{sub:4.destructuring}
	
	\emph{Destructuring} extrae valores de objetos y arrays con
	sintaxis declarativa concisa. \emph{Spread} (\texttt{...}) expande
	colecciones en distintos contextos; \emph{rest} recoge el
	«resto» de elementos en un parámetro. Véase el Listado~\ref{lst:4.4}.
	
	\begin{lstlisting}[style=js,caption={Destructuring, spread y rest},label=lst:4.4]
		// --- Destructuring de objeto ---
		const usuario = { nombre: "Ana", edad: 22, rol: "admin" };
		const { nombre, rol } = usuario;               // extrae dos campos
		const { edad: anios } = usuario;               // renombra: anios = 22
		const { pais = "EC" } = usuario;              // valor por defecto
		
		// --- Destructuring de array ---
		const colores = ["rojo", "verde", "azul"];
		const [primero, , tercero] = colores;          // omite el segundo
		
		// --- Rest en parametros ---
		function suma(...numeros) {
			return numeros.reduce((acc, n) => acc + n, 0);
		}
		console.log(suma(1, 2, 3, 4));                // 10
		
		// --- Spread para copiar/combinar ---
		const base   = { x: 1, y: 2 };
		const copia  = { ...base };                   // copia superficial
		const amplia = { ...base, z: 3 };             // extiende
		const nums   = [1, 2, 3];
		const mas    = [...nums, 4, 5];               // concatena
		
		// --- Destructuring en parametros ---
		function mostrar({ nombre, edad = 0 }) {
			console.log(`${nombre} tiene ${edad} anios`);
		}
		mostrar(usuario);
	\end{lstlisting}
	
	% ----- 4.3.5 Optional chaining y nullish coalescing -------------------
	\subsection{Optional chaining y nullish coalescing}
	\label{sub:4.optional}
	
	ES2020 introdujo dos operadores que simplifican el acceso seguro
	a propiedades que pueden no existir y el manejo de valores
	nulos o indefinidos. Véase el Listado~\ref{lst:4.5}.
	
	\begin{lstlisting}[style=js,numbers=none,caption={Optional chaining (\texttt{?.}) y
			nullish coalescing (\texttt{??})},label=lst:4.5]
		const resp = {
			data: {
				usuario: { nombre: "Luis", avatar: null }
			}
		};
		
		// Sin optional chaining (propenso a TypeError)
		// const nombre = resp.data.usuario.nombre; // ok
		// const ciudad = resp.data.usuario.ciudad.nombre; // TypeError!
		
		// Con optional chaining
		const nombre = resp?.data?.usuario?.nombre;       // "Luis"
		const ciudad = resp?.data?.usuario?.ciudad?.nombre; // undefined
		
		// Invocar metodo opcional
		const longitud = resp?.data?.usuario?.nombre?.length; // 4
		
		// Nullish coalescing: ?? retorna el lado derecho solo si
		// el izquierdo es null o undefined (NO si es 0 o "")
		const avatar = resp.data.usuario.avatar ?? "/img/default.png";
		const total  = 0 ?? 100;     // 0  (0 no es null ni undefined)
		const total2 = null ?? 100;  // 100
	\end{lstlisting}
	
	% ----- 4.3.6 Funciones de orden superior ------------------------------
	\subsection{Programación funcional: \texttt{map}, \texttt{filter} y
		\texttt{reduce}}
	\label{sub:4.functional}
	
	Los tres métodos de array más usados en JavaScript moderno
	implementan el paradigma funcional: transformar, filtrar y
	reducir datos sin mutar el array original. Véase el
	Listado~\ref{lst:4.6}.
	
	\begin{lstlisting}[style=js,caption={Funciones de orden superior sobre arrays},
		label=lst:4.6]
		const productos = [
		{ nombre: "Camiseta", precio: 15, activo: true  },
		{ nombre: "Pantalon", precio: 45, activo: true  },
		{ nombre: "Sombrero", precio:  8, activo: false },
		{ nombre: "Zapatos",  precio: 60, activo: true  },
		];
		
		// MAP: transforma cada elemento -> nuevo array
		const nombres = productos.map(p => p.nombre);
		// ["Camiseta", "Pantalon", "Sombrero", "Zapatos"]
		
		const conIva = productos.map(p => ({
			...p,
			precioFinal: +(p.precio * 1.15).toFixed(2)
		}));
		
		// FILTER: conserva solo los que cumplen la condicion
		const activos = productos.filter(p => p.activo);
		// 3 productos (sin Sombrero)
		
		// REDUCE: acumula un unico valor
		const totalActivos = productos
		.filter(p => p.activo)
		.reduce((suma, p) => suma + p.precio, 0);
		// 120
		
		// CADENA fluida: nombres de activos con precio < 50
		const baratos = productos
		.filter(p => p.activo && p.precio < 50)
		.map(p => p.nombre);
		// ["Camiseta", "Pantalon"]
	\end{lstlisting}
	
	% =====================================================================
	\section{Cierres (\emph{closures})}
	\label{sec:4.closures}
	
	Un \textbf{closure} es una función que «recuerda» el entorno
	léxico en el que fue creada, incluso después de que dicho entorno
	haya terminado de ejecutarse. Es uno de los conceptos más
	poderosos —y más incomprendidos— de JavaScript.
	
	\begin{cajadefinicion}{Definición formal de closure}
		Una \textbf{closure} es la combinación de una función y el
		entorno léxico en el que fue declarada. Ese entorno consiste en
		las variables locales que estaban en \emph{scope} en el momento
		de la creación de la función \cite{flanagan2020javascript}.
	\end{cajadefinicion}
	
	\begin{cajaejemplo}{Ejemplo 4.1 (RESUELTO) — Contador con closure}
		
		\textbf{Objetivo:} comprender cómo un closure mantiene estado
		privado sin variables globales.
		
		\begin{lstlisting}[style=js,caption={Contador encapsulado con closure},label=lst:4.7]
			"use strict";
			
			// La funcion crearContador devuelve un objeto
			// cuyas funciones comparten la variable privada 'cuenta'
			function crearContador(inicio = 0) {
				let cuenta = inicio;          // variable privada (no accesible desde fuera)
				
				return {
					incrementar()  { cuenta++;      return cuenta; },
					decrementar()  { cuenta--;      return cuenta; },
					reiniciar()    { cuenta = inicio; return cuenta; },
					valor()        { return cuenta; }
				};
			}
			
			const c1 = crearContador(10);
			const c2 = crearContador();     // cada instancia tiene su propio 'cuenta'
			
			console.log(c1.incrementar()); // 11
			console.log(c1.incrementar()); // 12
			console.log(c2.valor());       // 0  (independiente de c1)
			console.log(c1.reiniciar());   // 10
		\end{lstlisting}
		
		\textbf{Verificación:} intentar acceder a \texttt{c1.cuenta}
		devuelve \texttt{undefined}: la variable es verdaderamente
		privada. No existe ninguna API del navegador para acceder al
		cierre desde el exterior.
	\end{cajaejemplo}
	
	\begin{cajaejemplo}{Ejemplo 4.2 (RESUELTO) — Memoización con closure}
		
		\textbf{Objetivo:} cachear resultados de funciones costosas.
		
		\begin{lstlisting}[style=js,caption={Memoización genérica},label=lst:4.8]
			"use strict";
			
			function memoizar(fn) {
				const cache = new Map();                          // almacen privado
				
				return function (...args) {
					const clave = JSON.stringify(args);             // clave unica por argumentos
					if (cache.has(clave)) {
						console.log("HIT cache:", clave);
						return cache.get(clave);
					}
					const resultado = fn.apply(this, args);
					cache.set(clave, resultado);
					return resultado;
				};
			}
			
			// Funcion costosa de ejemplo
			function fibonacci(n) {
				if (n <= 1) return n;
				return fibonacci(n - 1) + fibonacci(n - 2);
			}
			
			const fibMemo = memoizar(fibonacci);
			console.log(fibMemo(35)); // calcula (tarda ~ms)
			console.log(fibMemo(35)); // HIT cache: instantaneo
		\end{lstlisting}
		
	\end{cajaejemplo}
	
	% =====================================================================
	\section{Prototipos y clases ES6}
	\label{sec:4.clases}
	
	JavaScript usa herencia basada en \emph{prototipos}, no en clases
	clásicas. La sintaxis \texttt{class} de ES6 es \emph{azúcar
		sintáctica} sobre el sistema de prototipos: lo hace más legible
	sin cambiar el modelo subyacente.
	
	\begin{cajahistoria}{El malentendido de las «clases» en JavaScript}
		Cuando Java y C++ dominaban la programación orientada a objetos,
		Brendan Eich eligió un modelo diferente inspirado en
		\textbf{Self} (Xerox PARC, 1986): los objetos heredan directamente
		de otros objetos a través de una cadena de prototipos, sin
		necesidad de clases. Esta decisión fue pragmática —diez días no
		dan para más— pero resultó enormemente poderosa.
		
		La comunidad pidió clases durante una década. ES6 (2015) las
		introdujo, pero como \emph{syntactic sugar}: internamente siguen
		siendo funciones constructoras y prototipos. El desarrollador
		que comprende esto entiende por qué
		\texttt{typeof} \texttt{MiClase} es \texttt{"function"} y por qué
		la cadena de prototipos sigue siendo accesible mediante
		\texttt{Object.getPrototypeOf()}.
	\end{cajahistoria}
	
	\begin{lstlisting}[style=js,caption={Clases ES6 con herencia y campos privados},
		label=lst:4.9]
		"use strict";
		
		// Clase base con campo privado (ES2022)
		class Figura {
			#color;                               // campo privado
			
			constructor(color = "negro") {
				this.#color = color;
			}
			
			get color() { return this.#color; }   // getter publico
			
			describir() {
				return `Figura de color ${this.#color}`;
			}
			
			// Metodo estatico: pertence a la clase, no a la instancia
			static esValido(obj) {
				return obj instanceof Figura;
			}
		}
		
		// Herencia con extends y super
		class Rectangulo extends Figura {
			#ancho;
			#alto;
			
			constructor(ancho, alto, color) {
				super(color);                       // llama al constructor padre
				this.#ancho = ancho;
				this.#alto  = alto;
			}
			
			get area()      { return this.#ancho * this.#alto; }
			get perimetro() { return 2 * (this.#ancho + this.#alto); }
			
			describir() {                         // sobrescribe (override)
				return `${super.describir()}, ${this.#ancho}x${this.#alto}cm`;
			}
		}
		
		const r = new Rectangulo(10, 5, "azul");
		console.log(r.describir());   // Figura de color azul, 10x5cm
		console.log(r.area);          // 50
		console.log(Figura.esValido(r)); // true
	\end{lstlisting}
	
	\begin{cajacompara}{Prototype vs.\ Class: ¿qué pasa bajo el capó?}
		\begin{lstlisting}[style=js,numbers=none]
			// Con clase ES6
			class Animal { hablar() { return "..."; } }
			
			// Es equivalente a (pre-ES6):
			function Animal() {}
			Animal.prototype.hablar = function() { return "..."; };
			
			// Verificacion
			console.log(typeof Animal);              // "function"
			console.log(Animal.prototype.hablar);   // [Function: hablar]
		\end{lstlisting}
	\end{cajacompara}
	
	% =====================================================================
	\section{Manipulación del DOM}
	\label{sec:4.dom}
	
	El \textbf{DOM} (\emph{Document Object Model}) es la
	representación en memoria de la página HTML, expuesta como un
	árbol de nodos que JavaScript puede leer y modificar. Cada
	elemento HTML es un objeto con propiedades y métodos.
	
	\begin{cajadefinicion}{El árbol del DOM}
		El DOM tiene una raíz (\texttt{document}), que contiene
		\texttt{document.documentElement} (\texttt{<html>}), y de ahí
		descienden \texttt{<head>} y \texttt{<body>}. Cada nodo puede
		ser de tipo \texttt{Element}, \texttt{Text} o \texttt{Comment}.
		La manipulación del DOM debe ser mínima y eficiente: cada
		modificación puede provocar \emph{reflow} y \emph{repaint}
		en el navegador.
	\end{cajadefinicion}
	
	\subsection{Selección de elementos}
	\label{sub:4.dom.select}
	
	\begin{lstlisting}[style=js,numbers=none,caption={Métodos de selección del DOM},
		label=lst:4.10]
		// Seleccion por ID (devuelve un elemento o null)
		const titulo = document.getElementById("titulo");
		
		// Seleccion con selector CSS (primero que coincide)
		const btn = document.querySelector(".btn-primario");
		
		// Seleccion multiple (devuelve NodeList estatica)
		const items = document.querySelectorAll("ul li");
		
		// Iteracion sobre NodeList
		items.forEach(item => console.log(item.textContent));
		// Con spread a Array para usar .map()
		const textos = [...items].map(li => li.textContent);
		
		// Relaciones entre nodos
		const padre  = btn.parentElement;
		const hijos  = padre.children;           // HTMLCollection
		const sig    = btn.nextElementSibling;
	\end{lstlisting}
	
	\subsection{Modificación del DOM}
	\label{sub:4.dom.modify}
	
	\begin{lstlisting}[style=js,numbers=none,caption={Crear, modificar y eliminar nodos},
		label=lst:4.11]
		// Crear nodo
		const parrafo = document.createElement("p");
		parrafo.textContent = "Hola, DOM!";       // texto plano (seguro vs XSS)
		parrafo.className   = "resaltado";
		parrafo.dataset.id  = "42";               // <p data-id="42">
		
		// Agregar al documento
		document.body.appendChild(parrafo);       // al final
		document.body.prepend(parrafo);           // al inicio
		// insertBefore(nuevo, referencia)
		
		// Modificar atributos
		parrafo.setAttribute("lang", "es");
		parrafo.removeAttribute("lang");
		
		// Modificar estilos en linea (preferir clases CSS)
		parrafo.style.color      = "red";
		parrafo.style.fontWeight = "bold";
		
		// Agregar/quitar clases
		parrafo.classList.add("activo");
		parrafo.classList.remove("resaltado");
		parrafo.classList.toggle("visible");      // alterna
		
		// Eliminar nodo
		parrafo.remove();
	\end{lstlisting}
	
	\subsection{Eventos y delegación}
	\label{sub:4.eventos}
	
	\begin{lstlisting}[style=js,caption={Eventos con addEventListener y delegación},
		label=lst:4.12]
		"use strict";
		
		// Patron basico
		const boton = document.querySelector("#enviar");
		boton.addEventListener("click", (e) => {
			e.preventDefault();                     // cancela accion por defecto
			console.log("clic en", e.target.id);
		});
		
		// Delegacion de eventos: un solo listener en el contenedor
		// en lugar de uno por cada hijo (eficiente para listas dinamicas)
		const lista = document.querySelector("#lista");
		lista.addEventListener("click", (e) => {
			const item = e.target.closest("li");    // busca el li mas cercano
			if (!item) return;                      // clic fuera de un item
			item.classList.toggle("completado");
			console.log("toggled:", item.dataset.id);
		});
		
		// Eventos de teclado
		document.addEventListener("keydown", (e) => {
			if (e.key === "Escape") cerrarModal();
			if (e.ctrlKey && e.key === "s") guardar(e);
		});
		
		// Eliminar listener (se necesita referencia a la funcion)
		function manejarScroll() { /* ... */ }
		window.addEventListener("scroll", manejarScroll);
		window.removeEventListener("scroll", manejarScroll);
	\end{lstlisting}
	
	\begin{cajaejemplo}{Ejemplo 4.3 (RESUELTO) — Lista de tareas dinámica}
		
		\textbf{Objetivo:} aplicar DOM, eventos, delegación y validación
		de entrada en una mini-aplicación completa.
		
		\begin{lstlisting}[style=html,numbers=none,caption={tareas.html},label=lst:4.13html]
			<!DOCTYPE html>
			<html lang="es">
			<head>
			<meta charset="UTF-8">
			<meta name="viewport" content="width=device-width, initial-scale=1">
			<title>Mis tareas</title>
			<style>
			body      { font-family: sans-serif; max-width: 600px; margin: 2rem auto; }
			ul        { list-style: none; padding: 0; }
			li        { display: flex; justify-content: space-between;
				padding: .5rem; border-bottom: 1px solid #eee; }
			li.hecha  { text-decoration: line-through; color: #999; }
			button    { cursor: pointer; }
			</style>
			</head>
			<body>
			<h1>Mis tareas</h1>
			<form id="form-tarea">
			<input id="texto-tarea" type="text" placeholder="Nueva tarea..."
			minlength="3" required>
			<button type="submit">Agregar</button>
			</form>
			<ul id="lista"></ul>
			<script src="tareas.js" type="module"></script>
			</body>
			</html>
		\end{lstlisting}
		
		\begin{lstlisting}[style=js,caption={tareas.js},label=lst:4.13js]
			"use strict";
			
			const form  = document.getElementById("form-tarea");
			const input = document.getElementById("texto-tarea");
			const lista = document.getElementById("lista");
			let   idAuto = 0;
			
			// Agregar tarea
			form.addEventListener("submit", (e) => {
				e.preventDefault();
				const texto = input.value.trim();
				if (texto.length < 3) {
					input.setCustomValidity("Minimo 3 caracteres");
					input.reportValidity();
					return;
				}
				input.setCustomValidity("");
				agregarTarea(texto);
				input.value = "";
				input.focus();
			});
			
			function agregarTarea(texto) {
				const li  = document.createElement("li");
				li.dataset.id = ++idAuto;
				li.textContent = texto;
				
				const btn = document.createElement("button");
				btn.textContent = "X";
				btn.setAttribute("aria-label", "Eliminar tarea");
				li.appendChild(btn);
				lista.appendChild(li);
			}
			
			// Delegacion: clic en lista para marcar/eliminar
			lista.addEventListener("click", (e) => {
				const li  = e.target.closest("li");
				if (!li) return;
				
				if (e.target.tagName === "BUTTON") {
					li.remove();                                // eliminar
				} else {
					li.classList.toggle("hecha");               // marcar completa
				}
			});
		\end{lstlisting}
		
		\textbf{Resultado esperado:}
		
		\fbox{\begin{minipage}{0.87\textwidth}
				\vspace{0.4em}
				{\Large\bfseries\color{uteqverde} Mis tareas}\\[0.5em]
				\fbox{\parbox{0.54\linewidth}{Nueva tarea\ldots}}
				\colorbox{uteqverde}{\textcolor{white}{\bfseries\ Agregar\ }}\\[0.6em]
				Estudiar JavaScript \hfill \fbox{X}\\[0.15em]\hrule
				{\color{gray} [completada] Hacer ejercicios HTML} \hfill \fbox{X}\\[0.15em]\hrule
				Practicar Flexbox \hfill \fbox{X}
				\vspace{0.4em}
		\end{minipage}}
		
		\medskip
		La segunda tarea está marcada como completada (tachada) al hacer
		clic sobre su texto. El botón «X» la elimina.
	\end{cajaejemplo}
	
	% =====================================================================
	\section{Programación asíncrona: callbacks, Promises y async/await}
	\label{sec:4.async}
	
	La asincronía es ineludible en el navegador: peticiones HTTP,
	lectura de archivos, temporizadores y eventos del usuario no deben
	bloquear la interfaz. JavaScript evolucionó a través de tres
	mecanismos sucesivos para manejarla.
	
	\subsection{Callbacks y el \emph{callback hell}}
	\label{sub:4.callbacks}
	
	El patrón más antiguo usa funciones de retorno
	(\emph{callbacks}) que se pasan como argumento. El problema
	surge cuando se encadenan: se forma una pirámide de código
	difícil de leer y mantener, conocida como \emph{callback hell}
	o \emph{pyramid of doom}. Véase el Listado~\ref{lst:4.14}.
	
	\begin{lstlisting}[style=js,numbers=none,caption={\emph{Callback hell} (patrón a evitar)},
		label=lst:4.14]
		// ANTI-PATRON: callback hell
		iniciarSesion(usuario, pass, (err, token) => {
			if (err) return manejarError(err);
			obtenerPerfil(token, (err, perfil) => {
				if (err) return manejarError(err);
				cargarPreferencias(perfil.id, (err, prefs) => {
					if (err) return manejarError(err);
					renderizar(perfil, prefs);        // profundidad maxima: ilegible
				});
			});
		});
	\end{lstlisting}
	
	\subsection{Promises}
	\label{sub:4.promises}
	
	Una \textbf{Promise} representa el resultado eventual de una
	operación asíncrona. Puede estar en tres estados:
	\texttt{pending}, \texttt{fulfilled} o \texttt{rejected}.
	Permite encadenar operaciones con \texttt{.then()} y capturar
	errores con \texttt{.catch()}, eliminando la pirámide. Véase el
	Listado~\ref{lst:4.15}.
	
	\begin{figure}[htbp]
		\centering
		\begin{tikzpicture}[
			state/.style={rectangle,draw=uteqverde,fill=uteqverde!8,thick,
				rounded corners=6pt,minimum width=2.8cm,minimum height=1cm,
				font=\small\bfseries,align=center},
			arr/.style={-{Latex[length=2.5mm]},thick},
			ok/.style={arr,color=uteqverde},
			fail/.style={arr,color=red!70!black}
			]
			\node[state,fill=gray!10,draw=gray!50] (pending) at (0,0) {pending};
			\node[state,fill=uteqverde!15] (fulfilled) at (4.5,1.2) {fulfilled};
			\node[state,fill=red!8,draw=red!60] (rejected) at (4.5,-1.2) {rejected};
			
			\draw[ok]   (pending.east) -- node[above,font=\scriptsize]{resolve(valor)} (fulfilled.west);
			\draw[fail] (pending.east) -- node[below,font=\scriptsize]{reject(error)} (rejected.west);
			
			\node[draw=uteqverde!40,fill=uteqverde!5,rounded corners,
			font=\scriptsize,align=center,minimum width=2.4cm]
			(then) at (8,1.2) {\texttt{.then(fn)}\\ejecuta \texttt{fn}};
			\node[draw=red!40,fill=red!5,rounded corners,
			font=\scriptsize,align=center,minimum width=2.4cm]
			(catch) at (8,-1.2) {\texttt{.catch(fn)}\\maneja error};
			\node[draw=gray!50,fill=gray!5,rounded corners,
			font=\scriptsize,align=center,minimum width=2.4cm]
			(finally) at (8,0) {\texttt{.finally(fn)}\\siempre};
			
			\draw[ok]   (fulfilled.east) -- (then.west);
			\draw[fail] (rejected.east)  -- (catch.west);
			\draw[ok,dashed]   (fulfilled.east) -- (finally.west);
			\draw[fail,dashed] (rejected.east)  -- (finally.west);
		\end{tikzpicture}
		\caption{Estados y transiciones de una Promise.}
		\label{fig:4.promise}
	\end{figure}
	
	\begin{lstlisting}[style=js,caption={Crear y encadenar Promises},label=lst:4.15]
		// Crear una Promise manualmente
		function esperar(ms) {
			return new Promise(resolve => setTimeout(resolve, ms));
		}
		
		// Encadenar con .then() / .catch() / .finally()
		esperar(500)
		.then(() => {
			console.log("Paso 1 completado");
			return esperar(300);              // devuelve otra Promise
		})
		.then(() => {
			console.log("Paso 2 completado");
			throw new Error("Fallo simulado");
		})
		.catch(err => {
			console.error("Error capturado:", err.message);
		})
		.finally(() => {
			console.log("Siempre se ejecuta");
		});
		
		// Promise.all: espera que TODAS se resuelvan
		const [datos, config] = await Promise.all([
		fetch("/api/datos").then(r => r.json()),
		fetch("/api/config").then(r => r.json())
		]);
		
		// Promise.allSettled: espera TODAS sin importar resultado
		const resultados = await Promise.allSettled([
		fetch("/api/a").then(r => r.json()),
		fetch("/api/b").then(r => r.json())
		]);
		resultados.forEach(r => {
			if (r.status === "fulfilled") console.log("OK:", r.value);
			else console.error("Fallo:", r.reason.message);
		});
		
		// Promise.race: la primera que resuelva (o rechace) gana
		const timeout = new Promise((_, reject) =>
		setTimeout(() => reject(new Error("Timeout")), 3000)
		);
		const dato = await Promise.race([fetch("/api/lento"), timeout]);
	\end{lstlisting}
	
	\subsection{\texttt{async}/\texttt{await}: el azúcar de las Promises}
	\label{sub:4.asyncawait}
	
	\texttt{async}/\texttt{await} es sintaxis que convierte el
	código asíncrono en algo que parece síncrono, eliminando el
	encadenamiento de \texttt{.then()}. Una función \texttt{async}
	\emph{siempre} devuelve una Promise. \texttt{await} suspende la
	ejecución de la función hasta que la Promise se resuelve o
	rechaza. Véase el Listado~\ref{lst:4.16}.
	
	\begin{lstlisting}[style=js,caption={\texttt{async}/\texttt{await} con manejo de errores},
		label=lst:4.16]
		"use strict";
		
		// Funcion auxiliar: verifica la respuesta HTTP
		async function respuestaSegura(res) {
			if (!res.ok) {
				const cuerpo = await res.text().catch(() => "");
				throw new Error(`HTTP ${res.status}: ${cuerpo}`);
			}
			return res.json();
		}
		
		// Funcion principal async
		async function obtenerUsuario(id) {
			try {
				const res     = await fetch(`https://jsonplaceholder.typicode.com/users/${id}`);
				const usuario = await respuestaSegura(res);
				return usuario;
			} catch (err) {
				console.error("Error al obtener usuario:", err.message);
				return null;                      // valor centinela
			}
		}
		
		// Llamada al nivel de modulo (top-level await en ES2022)
		const usuario = await obtenerUsuario(1);
		if (usuario) {
			console.log(`${usuario.name} <${usuario.email}>`);
		}
	\end{lstlisting}
	
	\subsection{Cancelar peticiones con AbortController}
	\label{sub:4.abort}
	
	\begin{lstlisting}[style=js,numbers=none,caption={Cancelar fetch con AbortController},
		label=lst:4.17]
		"use strict";
		
		let controlador = null;
		
		async function buscar(termino) {
			// Cancelar busqueda anterior si aun esta en curso
			if (controlador) controlador.abort();
			controlador = new AbortController();
			const { signal } = controlador;
			
			try {
				const res  = await fetch(
				`https://api.ejemplo.com/buscar?q=${encodeURIComponent(termino)}`,
				{ signal }
				);
				const data = await res.json();
				renderizarResultados(data);
			} catch (err) {
				if (err.name === "AbortError") {
					console.log("Busqueda cancelada (nueva consulta)");
				} else {
					console.error("Error de red:", err.message);
				}
			}
		}
		
		// Uso: buscar en tiempo real al escribir (debounced)
		document.querySelector("#buscador")
		.addEventListener("input", e => buscar(e.target.value));
	\end{lstlisting}
	
	\begin{cajaejemplo}{Ejemplo 4.4 (RESUELTO) — Panel de clima con fetch}
		
		\textbf{Objetivo:} consumir una API REST pública, mostrar datos
		en el DOM y manejar errores de red y de lógica de aplicación.
		
		\begin{lstlisting}[style=js,caption={clima.js -- modulo completo},label=lst:4.18]
			"use strict";
			
			const API_KEY = "DEMO_KEY";                   // reemplazar con clave real
			const URL_BASE = "https://api.open-meteo.com/v1/forecast";
			
			// Datos de ciudades (sustituye llamada a geocoding por brevedad)
			const CIUDADES = {
				quevedo:  { lat: -1.02, lon: -79.46 },
				guayaquil:{ lat: -2.18, lon: -79.89 },
				quito:    { lat: -0.22, lon: -78.51 }
			};
			
			async function obtenerClima(ciudad) {
				const coords = CIUDADES[ciudad.toLowerCase()];
				if (!coords) throw new Error(`Ciudad desconocida: ${ciudad}`);
				
				const params = new URLSearchParams({
					latitude:  coords.lat,
					longitude: coords.lon,
					current:   "temperature_2m,weathercode,wind_speed_10m",
					timezone:  "America/Guayaquil"
				});
				
				const res = await fetch(`${URL_BASE}?${params}`);
				if (!res.ok) throw new Error(`HTTP ${res.status}`);
				return res.json();
			}
			
			function renderClima(data, ciudad) {
				const { temperature_2m, wind_speed_10m } = data.current;
				document.getElementById("ciudad").textContent  = ciudad;
				document.getElementById("temp").textContent    = `${temperature_2m} degC`;
				document.getElementById("viento").textContent  = `${wind_speed_10m} km/h`;
			}
			
			// Manejador del boton
			document.getElementById("btn-clima").addEventListener("click", async () => {
				const ciudad = document.getElementById("input-ciudad").value.trim();
				const estado = document.getElementById("estado");
				
				if (!ciudad) { estado.textContent = "Ingrese una ciudad."; return; }
				
				estado.textContent = "Cargando...";
				try {
					const data = await obtenerClima(ciudad);
					renderClima(data, ciudad);
					estado.textContent = "";
				} catch (err) {
					estado.textContent = `Error: ${err.message}`;
				}
			});
		\end{lstlisting}
		
		\textbf{Cómo probarlo:} crear \texttt{clima.html} con un
		\texttt{<input id="input-ciudad">}, \texttt{<button id="btn-clima">},
		\texttt{<span id="ciudad">}, \texttt{<span id="temp">},
		\texttt{<span id="viento">} y \texttt{<span id="estado">}.
		Cargar en el navegador y escribir \texttt{quevedo}.
	\end{cajaejemplo}
	
	% =====================================================================
	\section{Módulos ES (\texttt{import}/\texttt{export})}
	\label{sec:4.modulos}
	
	Los módulos ES (ESM) son la solución nativa del lenguaje para
	organizar el código en ficheros independientes con un
	espacio de nombres propio. Cada módulo tiene su propio scope:
	no hay variables globales implícitas.
	
	\begin{cajahistoria}{Del global soup a los módulos ES}
		Durante más de quince años, el único mecanismo de modularización
		en JavaScript era concatenar scripts en el HTML o usar patrones
		como el \emph{Revealing Module Pattern} (2007, Christian Heilmann)
		basado en closures. La comunidad inventó soluciones como
		\textbf{CommonJS} (Node.js, 2009, con \texttt{require()}) y
		\textbf{AMD} (RequireJS, 2010, con \texttt{define()}). La guerra
		de formatos terminó con ES2015: los \textbf{módulos ES nativos}
		son el estándar oficial. Hoy todos los navegadores modernos y
		Node.js 12+ los soportan sin transpilación.
	\end{cajahistoria}
	
	\begin{lstlisting}[style=js,caption={Módulos ES: export e import},label=lst:4.19]
		// ============ archivo: matematica.js ============
		// Exportacion nombrada
		export const PI = 3.14159265;
		
		export function circunferencia(r) {
			return 2 * PI * r;
		}
		
		export function area(r) {
			return PI * r ** 2;
		}
		
		// Exportacion por defecto (solo una por modulo)
		export default class Calculadora {
			sumar(a, b)  { return a + b; }
			restar(a, b) { return a - b; }
		}
		
		// ============ archivo: main.js ============
		// Importacion nombrada
		import { PI, circunferencia, area } from "./matematica.js";
		
		// Importacion por defecto
		import Calculadora from "./matematica.js";
		
		// Importacion con alias
		import { circunferencia as circ } from "./matematica.js";
		
		// Importacion de todo el modulo como objeto
		import * as Mat from "./matematica.js";
		
		console.log(PI);                   // 3.14159265
		console.log(circunferencia(5));    // 31.416...
		const calc = new Calculadora();
		console.log(calc.sumar(3, 4));     // 7
		console.log(Mat.area(3));          // 28.274...
	\end{lstlisting}
	
	\begin{lstlisting}[style=html,numbers=none,caption={Usar módulos en HTML},
		label=lst:4.20html]
		<!-- type="module" activa el scope de modulo -->
		<!-- y habilita top-level await -->
		<script type="module" src="main.js"></script>
		<!-- El atributo defer es implicito en type="module" -->
	\end{lstlisting}
	
	\begin{cajaadvertencia}{CORS y módulos locales}
		Los módulos ES \textbf{no funcionan} cuando se abre el archivo
		HTML directamente con \texttt{file://}: el navegador bloquea la
		carga por política de origen cruzado (CORS). Se necesita un
		servidor HTTP. En IntelliJ IDEA basta con hacer clic derecho
		sobre el \texttt{.html} y elegir \textbf{Open in Browser}
		(IntelliJ inicia automáticamente un servidor local).
	\end{cajaadvertencia}
	
	\subsection{Importación dinámica}
	\label{sub:4.dynimport}
	
	La importación dinámica (\texttt{import()}) carga un módulo
	bajo demanda. Es clave para el \emph{code splitting} y la
	carga diferida (\emph{lazy loading}). Véase el
	Listado~\ref{lst:4.21}.
	
	\begin{lstlisting}[style=js,numbers=none,caption={Importación dinámica bajo demanda},
		label=lst:4.21]
		document.getElementById("btn-mapa")
		.addEventListener("click", async () => {
			// El modulo solo se carga cuando el usuario hace clic
			const { iniciarMapa } = await import("./mapa.js");
			iniciarMapa("#contenedor-mapa");
		});
	\end{lstlisting}
	
	% =====================================================================
	\section{Ejercicios paso a paso en IntelliJ IDEA}
	\label{sec:4.ejercicios}
	
	Los siguientes ejercicios incluyen instrucciones completas para
	reproducirlos en IntelliJ IDEA Ultimate con la licencia educativa
	gratuita de JetBrains.
	
	\begin{cajaejercicio}{Ejercicio 4.1 (PROPUESTO) --- Calculadora interactiva}
		
		\textbf{Enunciado.} Implementar una calculadora web con JavaScript
		puro que soporte las cuatro operaciones básicas y muestre el
		resultado en tiempo real.
		
		\textbf{Criterios de aceptación:}
		\begin{itemize}
			\item Usa \texttt{const}/\texttt{let} y funciones flecha.
			\item Valida división por cero con mensaje amigable.
			\item Tiene \texttt{"use strict"} y \texttt{type="module"}.
			\item El historial de operaciones se muestra con DOM.
			\item Limpia el historial con un botón.
		\end{itemize}
		
		\textbf{Pistas:} función \texttt{calcular(a, op, b)} devuelve
		el resultado o lanza un \texttt{Error}; un listener de
		\texttt{submit} en el formulario llama a esa función y actualiza
		el DOM.
	\end{cajaejercicio}
	
	\begin{cajaejercicio}{Ejercicio 4.2 (PROPUESTO) --- Buscador de Pokémon}
		
		\textbf{Enunciado.} Construir una página que consuma la PokeAPI
		pública (\url{https://pokeapi.co/api/v2/pokemon/\{nombre\}}) con
		\texttt{fetch} y \texttt{async/await}, mostrando imagen, tipos y
		estadísticas del Pokémon buscado.
		
		\textbf{Criterios de aceptación:}
		\begin{itemize}
			\item Maneja errores HTTP 404 con mensaje al usuario.
			\item Usa \texttt{AbortController} para cancelar búsquedas anteriores.
			\item Muestra un indicador de «cargando...» mientras espera la respuesta.
			\item El código está organizado en al menos dos módulos ES.
		\end{itemize}
	\end{cajaejercicio}
	
	\begin{cajaejercicio}{Ejercicio 4.3 (RESUELTO) --- Galería de usuarios con módulos}
		
		\textbf{Enunciado.} Crear una galería de usuarios consumiendo la
		API pública \url{https://randomuser.me/api/?results=12}. Separar
		la lógica en tres módulos: \texttt{api.js}, \texttt{ui.js} y
		\texttt{main.js}.
		
		\textbf{Pasos en IntelliJ IDEA Ultimate:}
		
		\bigskip
		\textbf{Paso 1 — Crear el proyecto.}
		
		En IntelliJ, abrir el menú \textbf{File $\rightarrow$ New $\rightarrow$
			Project\ldots}, seleccionar \textbf{Empty Project}, escribir el nombre
		\texttt{galeria-usuarios} y confirmar con \textbf{Create}.
		\emph{Por qué:} un proyecto vacío evita plantillas innecesarias para un
		ejercicio de JavaScript puro.
		
		\textbf{Paso 2 — Crear la estructura de carpetas.}
		
		En el panel \textbf{Project}, hacer clic derecho sobre la raíz y
		elegir \textbf{New $\rightarrow$ Directory}. Crear la carpeta
		\texttt{src}. Dentro de ella crear tres archivos JavaScript:
		\texttt{api.js}, \texttt{ui.js}, \texttt{main.js}, y en la raíz
		crear \texttt{index.html}.
		
		\textbf{Paso 3 — Escribir \texttt{api.js}.}
		
		\begin{lstlisting}[style=js,caption={src/api.js},label=lst:4.ej3.api]
			"use strict";
			
			const URL_USUARIOS = "https://randomuser.me/api/";
			
			export async function obtenerUsuarios(cantidad = 12) {
				const params = new URLSearchParams({ results: cantidad, nat: "es,us" });
				const res = await fetch(`${URL_USUARIOS}?${params}`);
				if (!res.ok) throw new Error(`HTTP ${res.status}`);
				const data = await res.json();
				return data.results;
			}
		\end{lstlisting}
		
		\emph{Qué hace:} exporta una función asíncrona que consulta la
		API y devuelve el array de usuarios. Lanza un error si la
		respuesta HTTP no es exitosa.
		
		\textbf{Paso 4 — Escribir \texttt{ui.js}.}
		
		\begin{lstlisting}[style=js,caption={src/ui.js},label=lst:4.ej3.ui]
			"use strict";
			
			export function crearTarjeta(usuario) {
				const { name, email, picture, location } = usuario;
				const tarjeta = document.createElement("article");
				tarjeta.className = "tarjeta";
				
				tarjeta.innerHTML = `
				<img src="${picture.large}" alt="Foto de ${name.first}">
				<h3>${name.first} ${name.last}</h3>
				<p>${email}</p>
				<p><small>${location.city}, ${location.country}</small></p>
				`;
				return tarjeta;
			}
			
			export function mostrarError(msg) {
				const contenedor = document.getElementById("galeria");
				contenedor.innerHTML =
				`<p class="error">Error: ${msg}. Intente de nuevo.</p>`;
			}
			
			export function mostrarCargando(activo) {
				document.getElementById("cargando").hidden = !activo;
			}
		\end{lstlisting}
		
		\emph{Qué hace:} encapsula toda la lógica de creación de nodos
		DOM. No mezcla lógica de red con presentación.
		
		\textbf{Paso 5 — Escribir \texttt{main.js}.}
		
		\begin{lstlisting}[style=js,caption={src/main.js},label=lst:4.ej3.main]
			"use strict";
			import { obtenerUsuarios }             from "./api.js";
			import { crearTarjeta, mostrarError, mostrarCargando } from "./ui.js";
			
			const galeria = document.getElementById("galeria");
			
			async function cargarGaleria() {
				mostrarCargando(true);
				galeria.innerHTML = "";
				try {
					const usuarios = await obtenerUsuarios(12);
					const fragmento = document.createDocumentFragment();
					usuarios.forEach(u => fragmento.appendChild(crearTarjeta(u)));
					galeria.appendChild(fragmento);         // un solo reflow
				} catch (err) {
					mostrarError(err.message);
				} finally {
					mostrarCargando(false);
				}
			}
			
			document.getElementById("btn-recargar")
			.addEventListener("click", cargarGaleria);
			
			cargarGaleria();                            // carga inicial
		\end{lstlisting}
		
		\emph{Por qué \texttt{DocumentFragment}:} agregar elementos
		uno a uno fuerza un \emph{reflow} por cada inserción. Un
		fragmento agrupa las inserciones y genera un único \emph{reflow},
		mejorando el rendimiento.
		
		\textbf{Paso 6 — Escribir \texttt{index.html}.}
		
		\begin{lstlisting}[style=html,numbers=none,caption={index.html},label=lst:4.ej3.html]
			<!DOCTYPE html>
			<html lang="es">
			<head>
			<meta charset="UTF-8">
			<meta name="viewport" content="width=device-width, initial-scale=1">
			<title>Galeria de Usuarios</title>
			<style>
			*, *::before, *::after { box-sizing: border-box; }
			body { font-family: system-ui, sans-serif; margin: 0; padding: 1rem;
				background: #f5f5f5; }
			h1   { color: #006633; text-align: center; }
			#galeria { display: grid;
				grid-template-columns: repeat(auto-fill, minmax(220px,1fr));
				gap: 1rem; margin-top: 1rem; }
			.tarjeta { background: white; border-radius: 8px; padding: 1rem;
				text-align: center; box-shadow: 0 2px 6px rgba(0,0,0,.1); }
			.tarjeta img { border-radius: 50%; width: 96px; height: 96px; }
			.tarjeta h3  { margin: .5rem 0 .25rem; font-size: 1rem; }
			.tarjeta p   { margin: 0; color: #555; font-size: .85rem; }
			.error { color: red; grid-column: 1/-1; text-align: center; }
			button { display: block; margin: 0 auto; padding: .6rem 1.4rem;
				background: #006633; color: white; border: none;
				border-radius: 6px; cursor: pointer; font-size: 1rem; }
			button:hover { background: #004d26; }
			#cargando { text-align: center; color: #888; margin-top: 2rem; }
			</style>
			</head>
			<body>
			<h1>Galeria de Usuarios</h1>
			<button id="btn-recargar">Recargar galeria</button>
			<p id="cargando" hidden>Cargando...</p>
			<section id="galeria" aria-live="polite"></section>
			<script type="module" src="src/main.js"></script>
			</body>
			</html>
		\end{lstlisting}
		
		\textbf{Paso 7 — Ejecutar y verificar.}
		
		En IntelliJ, hacer clic derecho sobre \texttt{index.html} y
		elegir \textbf{Open in} $\rightarrow$ \textbf{Chrome} (o el
		navegador preferido). IntelliJ levanta un servidor HTTP local
		en \texttt{localhost:63342} automáticamente. Abrir las DevTools
		(F12) $\rightarrow$ pestaña \textbf{Network} y verificar
		que la petición a \texttt{randomuser.me} devuelva 200 y un JSON
		con 12 usuarios. En la pestaña \textbf{Console} no debe aparecer
		ningún error.
		
		\emph{Verificación:} la galería debe mostrar 12 tarjetas con foto,
		nombre, correo y ciudad. Al hacer clic en «Recargar galería» se
		obtienen 12 usuarios distintos.
	\end{cajaejercicio}
	
	% =====================================================================
	\section{Buenas prácticas de JavaScript 2026}
	\label{sec:4.buenas}
	
	\begin{cajabuenas}{Doce reglas profesionales JavaScript 2026}
		\begin{enumerate}
			\item \textbf{\texttt{"use strict"}} en todo módulo o
			\textbf{\texttt{type="module"}} en el \texttt{<script>}
			(activa \emph{strict mode} implícitamente).
			\item \textbf{\texttt{const} por defecto}; \texttt{let} solo
			cuando la variable muta; \texttt{var} nunca.
			\item \textbf{\texttt{===} sobre \texttt{==}} siempre
			(comparación estricta evita coerciones silenciosas).
			\item \textbf{Manejo de errores} con \texttt{try/catch} en
			todas las funciones \texttt{async}; nunca silenciar un
			\texttt{catch} vacío.
			\item \textbf{Validar entradas} de funciones públicas antes de
			usarlas; lanzar \texttt{TypeError} con mensaje claro.
			\item \textbf{Funciones pequeñas}, una sola responsabilidad;
			máximo 20 líneas como heurística.
			\item \textbf{Sin variables globales}: usar módulos ES con
			\texttt{type="module"}.
			\item \textbf{Linter ESLint} con \texttt{eslint-config-airbnb}
			o \texttt{eslint-config-standard}; integrarlo en IntelliJ.
			\item \textbf{Tests} con Vitest o Jest; cobertura mínima 70\%
			en lógica de negocio.
			\item \textbf{No mezclar} \texttt{.then()} y \texttt{await}
			en la misma función: elegir un estilo y ser consistente.
			\item \textbf{Optional chaining} (\texttt{?.}) para acceso
			seguro a propiedades anidadas; nullish coalescing (\texttt{??})
			para valores por defecto precisos.
			\item \textbf{Tipado gradual} con JSDoc (\texttt{@param},
			\texttt{@returns}) o TypeScript en proyectos serios;
			IntelliJ IDEA proporciona completado de tipos con JSDoc.
		\end{enumerate}
	\end{cajabuenas}
	
	\begin{cajaadvertencia}{Anti-patrones frecuentes que se deben evitar}
		\begin{itemize}
			\item \textbf{\texttt{innerHTML} con datos del usuario}: riesgo
			de XSS. Usar \texttt{textContent} o
			\texttt{element.insertAdjacentText()} para texto sin
			etiquetas.
			\item \textbf{Callback hell}: usar \texttt{async/await}.
			\item \textbf{Polling} con \texttt{setInterval} para actualizar
			datos: usar WebSocket o Server-Sent Events.
			\item \textbf{Modificar el DOM en bucle}: usar
			\texttt{DocumentFragment}.
			\item \textbf{Olvidar cancelar peticiones}: usar
			\texttt{AbortController}.
			\item \textbf{Eval()}: nunca; riesgo de seguridad e impide
			optimizaciones del motor V8.
		\end{itemize}
	\end{cajaadvertencia}
	
	% =====================================================================
	\section{Diagrama de arquitectura: cliente con módulos ES}
	\label{sec:4.arquitectura}
	
	La Figura~\ref{fig:4.modulos} muestra la arquitectura típica
	de una aplicación de cliente moderna organizada con módulos ES.
	
	\begin{figure}[htbp]
		\centering
		\begin{tikzpicture}[
			mod/.style={rectangle,draw=uteqverde,fill=uteqverde!8,thick,
				rounded corners=4pt,minimum width=3cm,minimum height=1.1cm,
				font=\small\bfseries,align=center},
			ext/.style={rectangle,draw=uteqdorado!70!black,fill=uteqdorado!10,thick,
				rounded corners=4pt,minimum width=2.8cm,minimum height=1cm,
				font=\small,align=center},
			arr/.style={-{Latex[length=2mm]},thick,gray!70!black},
			imp/.style={arr,dashed,uteqverde!60!black},
			lbl/.style={font=\scriptsize,color=gray!70!black}
			]
			% HTML
			\node[ext,minimum width=2.2cm,draw=blue!50,fill=blue!5]
			(html) at (0,5.5) {\texttt{index.html}\\
				\scriptsize type="module"};
			
			% Modulos
			\node[mod] (main) at (0,3.5)   {\texttt{main.js}\\
				\scriptsize orquesta};
			\node[mod] (api)  at (-3.5,1.5) {\texttt{api.js}\\
				\scriptsize \texttt{fetch}/HTTP};
			\node[mod] (ui)   at (3.5,1.5)  {\texttt{ui.js}\\
				\scriptsize DOM/render};
			\node[mod] (util) at (0,1.5)   {\texttt{utils.js}\\
				\scriptsize helpers puros};
			
			% APIs externas
			\node[ext] (rest) at (-3.5,-0.5) {API REST\\
				\scriptsize JSON};
			\node[ext] (dom)  at (3.5,-0.5)  {DOM del\\navegador};
			
			% Flechas
			\draw[imp] (html) -- node[right,lbl]{carga} (main);
			\draw[imp] (main) -- node[above left,lbl,pos=0.4]{import} (api);
			\draw[imp] (main) -- node[above right,lbl,pos=0.4]{import} (ui);
			\draw[imp] (main) -- node[right,lbl]{import} (util);
			\draw[imp] (api)  -- node[right,lbl]{import} (util);
			\draw[arr] (api)  -- node[right,lbl]{\texttt{fetch}} (rest);
			\draw[arr] (ui)   -- node[right,lbl,align=left]{querySelector\\\\createElement} (dom);
		\end{tikzpicture}
		\caption{Arquitectura de una SPA mínima con módulos ES nativos.}
		\label{fig:4.modulos}
	\end{figure}
	
	% =====================================================================
	\section*{Resumen de la semana}
	\addcontentsline{toc}{section}{Resumen de la semana}
	
	\begin{cajadefinicion}{Lo esencial de la semana 4}
		\begin{itemize}
			\item JavaScript es \textbf{single-threaded} con un \emph{event
				loop} que permite asincronía sin bloquear la interfaz.
			\item \textbf{\texttt{const}/\texttt{let}}, funciones flecha,
			template literals, \emph{destructuring} y \emph{spread} son
			los pilares del JS moderno; \texttt{var} está obsoleto.
			\item Las \textbf{funciones de orden superior}
			(\texttt{.map()}, \texttt{.filter()}, \texttt{.reduce()})
			transforman datos de forma funcional e inmutable.
			\item Los \textbf{closures} permiten encapsular estado privado
			sin necesidad de clases.
			\item Las \textbf{clases ES6} son azúcar sobre prototipos;
			los campos privados (\texttt{\#}) datan de ES2022.
			\item La \textbf{manipulación del DOM} debe ser mínima y usar
			delegación de eventos para listas dinámicas.
			\item \textbf{\texttt{async/await}} simplifica el código
			asíncrono; siempre envolver con \texttt{try/catch}.
			\item Los \textbf{módulos ES} eliminan el scope global y son
			el estándar nativo; \texttt{AbortController} cancela peticiones.
		\end{itemize}
	\end{cajadefinicion}
	
	% =====================================================================
	% Bibliografía
	% =====================================================================
	\bibliographystyle{plain}
	\bibliography{D:/Repositorios/Asignaturas/AplicacionesWeb/Texto/LibroWebApp}
	
\end{document}